% style: article — clean single-column tutorial (default)
\usepackage{mathtools}
% Bold math. Pandoc's XeLaTeX template loads unicode-math, which breaks bm's \bm and
% amsmath's \boldsymbol on Greek letters ("Improper alphabetic constant") and makes
% \mathbf drop Greek glyphs. Route all three through unicode-math's math-font alphabets.
\makeatletter
\@ifpackageloaded{unicode-math}{%
  \unimathsetup{mathbf=sym}%
  \DeclareRobustCommand{\bm}[1]{\symbfit{#1}}%
  \DeclareRobustCommand{\boldsymbol}[1]{\symbfit{#1}}%
}{\usepackage{bm}}
\makeatother
\usepackage{float}
\floatplacement{figure}{H}
\usepackage[font=small,labelfont=bf]{caption}
\usepackage{booktabs}
\setlength{\parskip}{0.4em}
% Side notes: every `>` blockquote becomes a tutornote box (see html2latex.lua).
\usepackage[most]{tcolorbox}
\newtcolorbox{tutornote}{enhanced, breakable, colback=black!4, colframe=black!40,
  boxrule=0pt, leftrule=2.5pt, arc=0pt, outer arc=0pt, left=8pt, right=8pt, top=4pt, bottom=4pt,
  before skip=10pt, after skip=10pt, fontupper=\small}
% Korean/CJK next to Latin: no extra gap where the source has none ("map은", not "map 은").
% Spaces written in the source are kept (xeCJK's default xCJKecglue=false).
\makeatletter
% The en dash (–) is set in the Latin font: xeCJK would otherwise use the CJK font, and many
% Korean fonts (e.g. 조선굵은명조) have no en dash.
\@ifpackageloaded{xeCJK}{\xeCJKsetup{CJKecglue={}}\xeCJKDeclareCharClass{Default}{"2013,"2014,"2018,"2019,"201C,"201D}}{}
\makeatother
